\documentclass[11pt, oneside]{article}   	% use "amsart" instead of "article" for AMSLaTeX format
\usepackage{geometry}                		% See geometry.pdf to learn the layout options. There are lots.
\geometry{letterpaper}                   		% ... or a4paper or a5paper or ... 

\usepackage[parfill]{parskip}    		% Activate to begin paragraphs with an empty line rather than an indent
\usepackage{graphicx}				% Use pdf, png, jpg, or eps§ with pdflatex; use eps in DVI mode
								% TeX will automatically convert eps --> pdf in pdflatex		
\usepackage{amssymb}
\usepackage{cite}

% numbered examples
\usepackage{gb4e}
\usepackage{enumitem}
\usepackage{cancel}
\usepackage{amsmath}

\usepackage{mhchem}
\usepackage{lewis}
%\usepackage{urlbst}
\usepackage{url}
\usepackage{hyperref}


% Scientific Laws
\newtheorem{definition}{Definition}
\newtheorem{law}{Law}
\newtheorem{theory}{Theory}
\newtheorem{hint}{Hint}

\title{Module 12: The Gas Phase}
\author{Mark Peever\\ \texttt{mpeever@gmail.com}}
\date{March 6 -- 13, 2026}

\begin{document}
\maketitle

\section{Overview}
\begin{enumerate}
\item \textbf{Boyle's Law} relates pressure and volume.
\item \textbf{Charles's Law} relates temperature and volume.
\item \textbf{The Ideal Gas Law} relates pressure, temperature, volume, and number of moles.
\item \textbf{Dalton's Law of Partial Pressures} tells us how to combine pressures of gases to get totals.
\end{enumerate}

\section{Pressure}

\begin{definition}[Pressure]\label{def:pressure}
The force per unit area exerted on an object.\\
\end{definition}


\begin{itemize}
\item pressure is force per unit area: $P = \frac{F}{A}$
\item pressure is measured in Pascals: $1 Pa = \frac{1 N}{1 m^{2}}$
\item in English units, we measure pressure in \emph{pounds per square inch} (psi)
\item pressure is the result of the particles (atoms, molecules) of a substance bounding off surfaces
\item we can express pressure in atmospheres: $1.000 atm = 101.3 kPa$
\end{itemize}

\section{Boyle's Law}

\begin{law}[Boyle's Law]\label{law:boyles}
For any given temperature, the product of a gas's pressure and volume is constant.\\
$$P V = constant$$
\end{law}

\begin{itemize}
\item $P V = constant$ at any given temperature
\item so if you decrease volume, you increase pressure
\item this works clearly with kinetic theory: the less space particles have, the more they'll bump into things
\item we can express Boyle's Law like this: $P_1 V_1 = P_2 V_2$
\end{itemize}

\section{Charles's Law}

\begin{law}[Charles's Law]\label{law:charles}
At a constant pressure, the temperature and volume of a gas are linearly proportional.\\
$$\frac{V}{T} = constant$$
\end{law}

\begin{itemize}
\item Boyle's Law says pressure and volume are inversely proportional (as one goes up, the other goes down) (see Law \ref{law:boyles}, p. \pageref{law:boyles})
\item Charles's Law says temperature and volume are linearly proportional (as one goes up, the other goes up) 
\item this also makes sense with kinetic theory: the more we restrain objects from moving, the more they bump into each other
\item we can express Charles's Law like this: $\frac{V_1}{T_1} = \frac{V_2}{T_2}$
\end{itemize}

\section{Combined Gas Law}

\begin{law}[Combined Gas Law]\label{law:gas:combined}
$$\frac{P V}{T} = constant$$
\end{law}

\begin{itemize}
\item we can combine Charles's Law and Boyle's Law to get a Combined Gas Law:\\
$$\frac{P_1 V_1}{T_1} = \frac{P_2 V_2}{T_2}$$
\item we need to measure our temperatures in Kelvin to use this law\footnote{Yes, we can use Rankine as well. The point is we can't have an arbitrary zero, and both Fahrenheit and Celsius scales have arbitrary zeroes.}
\end{itemize}

\section{Ideal Gases}

\begin{definition}[Ideal Gas]\label{def:ideal-gas}
An ideal gas has the following properties:\\
\begin{enumerate}
\item The molecules or atoms that make up the gas are very small compared to the total volume available to the gas.
\item The molecules of atoms that make up the gas must be so far apart from one another that there is no attraction or repulsion between them.
\item The collisions that occur between the gas molecules or atoms must be elastic. When they collide, no energy can be lost in the collision.
Also, no energy can be lost when they collide with the walls of the container.
\end{enumerate}
\end{definition}

\begin{definition}[Standard Temperature and Pressure]\label{defn:stp}
A temperature of 273k and a pressure of 1.00 atm.
\end{definition}

\begin{itemize}
\item gases behave like ideal gases (see Definition \ref{def:ideal-gas}, p. \pageref{def:ideal-gas}) at high temperatures and low pressures
\item we define the ``standard" temperature and pressure (STP) as 273k and 1.00 atm (see Definition \ref{defn:stp}, p, \pageref{defn:stp})
\item so gases are pretty close to ideal at STP
\end{itemize}

\section{Partial Pressures}

\begin{law}[Dalton's Law of Partial Pressures]\label{law:dalton:partial-pressures}
When two or more ideal gases are mixed together, the total pressure of the mixture is equal to the sum of the pressures of each individual gas.\\
$$P_{mixture} = \sum^{n}_{i=1} P_{i} $$
\end{law}

\begin{hint}[Ideal Gas Identity]\label{hint:ideal-gas-identity}
The pressure of an ideal gas does not depend on the identity of the gas.
It depends only on the quantity of the gas.
\end{hint}

\begin{itemize}
\item the pressure of a mixture of gases is the sum of the pressure of each gas
\item in other words\ldots if we mix two gases with two different pressures, the mixture's pressure will be the sum of their pressures
\end{itemize}


\section{Vapor Pressure}

\begin{definition}[Vapor Pressure]\label{defn:vapor-pressure}
The pressure exerted by the vapor which sits on top of any liquid.
\end{definition}

\begin{definition}[Boiling Point]\label{defn:boiling-point}
The temperature at which the vapor pressure of a liquid is equal to the normal atmospheric pressure.
\end{definition}

\begin{itemize}
\item any given liquid has a small amount of vapor sitting above it, which results from evaporation of the liquid
\item because the vapor (gas) exerts pressure, there is a \emph{vapor pressure} associated with that liquid (see Definition \ref{defn:vapor-pressure}, p. \pageref{defn:vapor-pressure})
\item so the layer of evaporated liquid is pushing down on the liquid
\item the vapor pressure increases with temperature, as our combined gas law (Law \ref{law:gas:combined}, p. \pageref{law:gas:combined}) predicts
\item if we increase the temperature until the vapor pressure reaches atmospheric pressure, then liquid is at the boiling point (Definition \ref{defn:boiling-point}, p. \pageref{defn:boiling-point})
\end{itemize}

\section{Concentration: Molar Fraction}
\begin{itemize}
\item we can think of gas mixtures like solutions
\item just like solutions, we can think in terms of \emph{concentration}
\item with gases, we think in terms of \textbf{mole fraction}, which is just $\frac{moles of one substance}{moles of the whole mixture}$
\item so if we mix $1.000 mol$ of \ce{O2} into $3.000 mol$ of \ce{N2}, then the mole fraction of \ce{O2} is $\frac{1.000 mol}{1.000 mol + 3.000 mol} = 0.2500$
\item so in this case, the partial pressure of \ce{O2} would be $\frac{1}{4}$ of the pressure of the mixture
\end{itemize}


\section{Ideal Gas Law}

\begin{law}[Ideal Gas Law]\label{law:ideal-gas}
$$P V = n R T$$ \\
$$ R = 0.0821 \frac{L \cdot atm}{mole \cdot K}$$
\end{law}

\begin{itemize}
\item for an ideal gas (see Definition \ref{def:ideal-gas}, p. \pageref{def:ideal-gas}), we can related temperature, pressure, and volume by the Ideal Gas Law (Law \ref{law:ideal-gas}, p. \pageref{law:ideal-gas})
\item in that law:
\begin{itemize}
\item $n$ refers to the number of moles of our ideal gas
\item $R$ is a constant: $ R = 0.0821 \frac{L \cdot atm}{mole \cdot K}$
\item $P$ is pressure
\item $V$ is volume
\item $T$ is temperature
\end{itemize}
\item it shouldn't surprise us that we can use this in stoichiometery, because it lets us find $n$, the number of moles of an ideal gas!!!!
\end{itemize}

\section{Homework}
Review Problems: p. 415 \# 1--10 (not to be turned in)\\
Practice Problems: p. 416 \# 1--10 (due 2026-03-27)\\

\nocite{wile-chem-2}
\bibliography{../Chemistry}
\bibliographystyle{apalike}

\clearpage
\subsection{Answers to Practice Problems}
Answers to Practice Problems p. 416 \# 1--10:
\begin{enumerate}
\item $1.27 L$

\item $43 mL$

\item $137.7 atm$

\item $760 torr$

\item 
\begin{enumerate}
\item \ce{SO2} : $0.350$
\item \ce{NO} : $0.603$
\item \ce{SO3}: $0.046$
\end{enumerate}

\item 
\begin{enumerate}
\item \ce{SO2} : $0.39 atm$
\item \ce{NO} : $0.66 atm$
\item \ce{SO3}: $0.051 atm$
\end{enumerate}

\item 
\begin{enumerate}
\item $P_{total} = 7.50 atm$
\item mole fraction of \ce{N2} : $0.667$
\item mole fraction of \ce{O2} : $0.0.267$
\item mole fraction of \ce{Ar}: $0.067$
\end{enumerate}

\item $5.19 L$

\item $0.621 g$

\item $2790 L$

\end{enumerate}
\end{document}  